\question[3]
Beim Besuch vieler Websites erscheint zunächst ein Cookie-Banner, das 
um Zustimmung zur Verwendung von Cookies bittet.

\begin{enumerate}
    \item Erkläre, was Cookies sind und welche personenbezogenen 
    Daten mit ihrer Hilfe über Nutzerinnen und Nutzer gesammelt werden 
    können.
    \item Erläutere den Unterschied zwischen notwendigen Cookies und 
    Tracking-Cookies.
\end{enumerate}

\begin{solutionbox}{8cm}

\textbf{1.} Cookies sind kleine Textdateien, die eine Website beim 
Besuch im Browser des Nutzers speichert und bei einem erneuten Besuch 
wieder ausliest. Über Cookies können unter anderem gespeichert werden: 
Login-Zustände, Spracheinstellungen, der Inhalt eines Warenkorbs, 
besuchte Seiten und Verweildauer, sowie eine eindeutige Kennung (ID), 
mit der ein Nutzer bei späteren Besuchen wiedererkannt werden kann. In 
Kombination lassen sich daraus Nutzungsprofile erstellen, etwa welche 
Interessen jemand hat oder welche Seiten regelmäßig besucht werden.

\textbf{2.} Notwendige Cookies sind für die Grundfunktion einer 
Website erforderlich, z.\,B.\ um einen Nutzer während einer Sitzung 
eingeloggt zu halten oder den Warenkorb zu speichern; sie dienen nicht 
der Datensammlung zu anderen Zwecken. Tracking-Cookies dagegen werden 
gezielt eingesetzt, um das Verhalten von Nutzern über mehrere Websites 
hinweg zu verfolgen, meist zu Zwecken der Werbung oder Marktanalyse. 
Sie stammen häufig von Drittanbietern (z.\,B.\ Werbenetzwerken), die 
auf vielen verschiedenen Websites eingebunden sind und so ein 
umfassenderes Bewegungsprofil im Web erstellen können.

\end{solutionbox}


